\documentclass[conference]{IEEEtran}
\usepackage[T1]{fontenc}
\usepackage[utf8]{inputenc}
\usepackage{cite}
\usepackage{url}
\usepackage{booktabs}
\usepackage{array}
\usepackage{graphicx}
\usepackage{amsmath}
\usepackage{hyperref}
\hypersetup{hidelinks}

\title{Traffic Optimisation and Route Redistribution System: A Deep Reinforcement Learning Approach for Network-Wide Congestion Mitigation in Malta}

\author{
\IEEEauthorblockN{Kyle Austin}\\
\IEEEauthorblockA{
B.Sc. (Hons.) Applied Data Science / Software Development\\
MCAST Institute of Information and Communication Technology\\
Malta\\
Kyle.austin.h63676@mcast.edu.mt}
}

\begin{document}
\maketitle

%% -----------------------------------------------------------
%%  ABSTRACT
%% -----------------------------------------------------------
\begin{abstract}
Malta is ranked the world's second most congested nation, with drivers losing
approximately 94 hours annually to rush-hour delays on an island measuring only
27\,km $\times$ 14\,km. Unlike large metropolitan environments studied in the
existing literature, Malta's physically bounded road network cannot be expanded,
making intelligent traffic management the only scalable intervention.
This paper presents a prototype Traffic Optimisation and Route Redistribution
System (TORRS) that applies Multi-Agent Advantage Actor-Critic (MA2C), a
decentralised deep reinforcement learning algorithm, to jointly control traffic
signals and vehicle routing across 190 cooperative agents in a SUMO
microsimulation of the Maltese road network derived from OpenStreetMap data.
The system is evaluated against a greedy signal controller baseline over
3600-second simulation episodes.
After 14 training episodes the MA2C agent achieves a 37.3\% reduction in
average network waiting time (from 9.87\,s to 6.19\,s per vehicle), demonstrating
a clear learning signal despite a severely limited training budget.
The results establish a working prototype for network-wide RL-based traffic
management in a small-island context and identify computational scalability as
the primary challenge for future work.
\end{abstract}

\begin{IEEEkeywords}
Multi-Agent Reinforcement Learning, Traffic Signal Control, Vehicle Routing,
SUMO Microsimulation, Malta Traffic Congestion
\end{IEEEkeywords}

%% -----------------------------------------------------------
%%  I. INTRODUCTION
%% -----------------------------------------------------------
\section{Introduction}

Urban congestion remains one of the most persistent challenges in intelligent
transportation systems (ITS), imposing large economic losses, excess fuel
consumption, and degraded travel-time reliability on millions of daily commuters.
In most metropolitan contexts, road-capacity expansion remains a viable, if
expensive, lever for relief. However, in small island states with physically
bounded road networks, capacity expansion is not financially, environmentally,
or geographically feasible, making the intelligent redistribution of existing
traffic flows the only scalable intervention available.

Malta presents one of the most extreme instances of this problem. Ranked the
world's second most congested nation in 2026 by the TomTom Traffic Index, the
island hosts more than 454{,}000 registered vehicles for a population of
approximately 574{,}250 residents (roughly 79 vehicles per 100 people), one of
the highest ratios in the European Union~\cite{tomtom2026,nso2024}. Drivers in
Valletta, the capital, lose approximately 94 hours per year to rush-hour delays,
and 36 new vehicles are added to the national fleet every day. Despite the
introduction of free public transport, congestion has continued to rise, indicating
that supply-side and routing-based interventions must play a central role in any
solution~\cite{nso2024}.

The Traffic Optimisation and Route Redistribution System (TORRS) proposed in
this study addresses this gap. Rather than optimising the travel time of a single
vehicle, TORRS aims to improve collective network performance by combining joint
signal control and vehicle rerouting through a cooperative multi-agent
reinforcement learning framework. The central hypothesis is that a learning-based
controller can demonstrate measurable reductions in average vehicle waiting time
and travel time compared with a rule-based greedy baseline, and that performance
improves monotonically with training experience.

The study is guided by two research questions:

\begin{itemize}
    \item \textbf{RQ1:} How does the MA2C learning-based controller compare with
    a greedy rule-based controller in terms of average vehicle waiting time and
    travel time in a SUMO microsimulation of the Maltese road network?
    \item \textbf{RQ2:} Does the MA2C agent demonstrate a measurable learning
    trend over training episodes, evidencing that the framework can improve with
    additional training?
\end{itemize}

The remainder of the paper is structured as follows. Section~II reviews the
relevant literature. Section~III describes the research methodology and
simulation pipeline. Section~IV presents the experimental findings. Section~V
discusses implications, and Section~VI concludes with limitations and future
directions.

%% -----------------------------------------------------------
%%  II. LITERATURE REVIEW
%% -----------------------------------------------------------
\section{Literature Review}

\subsection{Overview of ITS and Traffic Optimisation}

Intelligent transportation systems integrate users, vehicles, and infrastructure
with communication and information technologies to improve safety, efficiency,
and environmental performance. A systematic review by Zulkarnain~\textit{et~al.}
analysed more than 20{,}000 ITS publications and identified reinforcement learning,
deep learning, and graph-based neural networks as the dominant emerging techniques
for traffic flow prediction and routing control~\cite{zulkarnain2021}. Despite
the rapid growth of ITS research, the same review highlights a persistent gap:
the overwhelming majority of studies focus on large metropolitan areas in North
America, East Asia, and Northern Europe. Small island states with unique
structural constraints (high vehicle density, bounded geography, and limited
parallel corridors) remain almost entirely absent from the literature.

\subsection{Dynamic Rerouting with Fog-Cloud Reinforcement Learning}

Du~\textit{et~al.} propose GAQ-EBkSP, a two-stage dynamic rerouting framework
that combines a Graph Attention Q-network with entropy-balanced $k$-shortest path
selection under a fog-cloud architecture~\cite{du2023}. Fog nodes aggregate local
road information and reduce the action space, while cloud servers coordinate
global optimisation. Evaluated in SUMO, the system achieves approximately
14--54\% improvements in mean travel speed relative to rule-based baselines,
particularly at low rerouting ratios. The fog-cloud partitioning maps naturally
onto Malta's geographic traffic zones (e.g., Valletta gateway, St~Julian's,
Msida corridor), making this framework directly relevant as a conceptual
reference for TORRS.

\subsection{Adaptive Rerouting and Secondary Bottlenecks}

Gregucci~\textit{et~al.} evaluate adaptive rerouting strategies on simulated
grid networks and demonstrate that naive shortest-path diversion consistently
overloads apparently optimal detours, shifting congestion rather than
eliminating it~\cite{gregucci2023}. Their topology-aware strategy monitors
segment utilisation across the whole network to prevent secondary bottleneck
formation. This finding is directly relevant to Malta, where the limited number
of alternative corridors between major zones (notably the approaches to Valletta
and the northern residential belt) means that poorly designed rerouting could
worsen overall network performance.

\subsection{MA2C for Joint Signal and Routing Control}

Chu~\textit{et~al.} introduce MA2C, a decentralised multi-agent advantage
actor-critic algorithm with LSTM memory and a neighbourhood fingerprinting
mechanism that partially addresses non-stationarity in cooperative multi-agent
settings~\cite{chu2020}. The reference implementation evaluates joint traffic
signal and routing control on the Sioux Falls benchmark network, achieving
convergence over approximately 694 training episodes. This implementation
serves as the algorithmic foundation for TORRS.

\subsection{Synthesis and Research Gap}

The reviewed literature establishes that (i)~learning-based dynamic rerouting
can outperform rule-based methods in urban simulations, (ii)~network-wide
awareness and secondary-bottleneck prevention are critical design requirements,
and (iii)~no existing study applies these approaches to a small island road
network. TORRS contributes to the literature by targeting this unexplored context
and demonstrating the feasibility of the MA2C framework on a real-world
island road network.

%% -----------------------------------------------------------
%%  III. RESEARCH METHODOLOGY
%% -----------------------------------------------------------
\section{Research Methodology}

\subsection{Research Pipeline}

The research follows a structured pipeline consisting of five stages:
(1) network construction and validation, (2) demand modelling,
(3) baseline establishment, (4) MA2C training, and (5) comparative
evaluation and statistical analysis.

\subsection{Simulation Environment}

The Maltese road network was imported into SUMO using \texttt{netconvert} applied
to an OpenStreetMap export~\cite{lopez2018}, and validated against Transport
Malta infrastructure data. The simulation runs at a 5-second control interval
over 3600-second episodes, with 442 lane-area detectors monitoring all
traffic-light-controlled incoming lanes. A TraCI Python interface exposes the
environment to the RL agent.

\subsection{Demand Modelling}

A base weekday peak-hour demand scenario was calibrated against NSO vehicle
registration statistics~\cite{nso2024}. Vehicles are generated from an
origin-destination matrix derived from available open datasets, with random seed
variation used to produce independent episode realisations for training and
evaluation. Extension to seasonal tourist-demand scenarios is identified as
future work.

\subsection{Agent Architecture}

MA2C assigns one agent to each of 17 signalised junctions and one agent to each
of 12 routing corridors, giving 190 cooperative agents in total (including
sub-agents per signal phase). Each agent maintains an LSTM-based actor-critic
network with a local observation vector derived from queue lengths, vehicle
waiting times, and neighbour fingerprints. Signal agents are penalised for
vehicle queue lengths and waiting times; routing agents are rewarded for vehicles
completing trips on their assigned corridor.

Key hyperparameters are listed in Table~\ref{tab:hyperparams}.

\begin{table}[h]
\caption{MA2C Hyperparameters}
\label{tab:hyperparams}
\centering
\renewcommand{\arraystretch}{1.2}
\begin{tabular}{lc}
\toprule
\textbf{Parameter} & \textbf{Value} \\
\midrule
Learning rate                          & $2.5 \times 10^{-4}$ \\
Discount factor $\gamma$               & $0.99$ \\
$n$-step return                        & $144$ \\
Cooperative discount (signal-signal)   & $0.3$ \\
Episode length                         & $3600$\,s \\
Control interval                       & $5$\,s \\
Training steps                         & $10{,}080$ (14 episodes) \\
\bottomrule
\end{tabular}
\end{table}

\subsection{Baseline Controller}

The greedy baseline (\texttt{MaltaController}) assigns each junction the signal
phase that maximises the combined vehicle count on the highest-flow pair of
incoming lanes at the current timestep. Route agents are fixed at action~0
(no rerouting), replicating static route assignment. The baseline was evaluated
over two 3600-second episodes with random seeds 10{,}000 and 20{,}000.

\subsection{Evaluation Protocol}

After training, the MA2C model was evaluated deterministically (argmax policy)
over ten random seeds (10{,}000--100{,}000) using SUMO's
\texttt{tripinfo-output} to record per-vehicle travel and waiting times,
yielding 9{,}817 complete vehicle trips for analysis. Primary metrics are
average waiting time and average travel time per completed trip.

%% -----------------------------------------------------------
%%  IV. FINDINGS
%% -----------------------------------------------------------
\section{Findings}

\subsection{RQ2: Learning Trend}

Table~\ref{tab:training} reports per-episode metrics from the training run.
Average waiting time decreases from 9.87\,s in episode~1 to 6.19\,s in
episode~14, a \textbf{37.3\% reduction} over 14 episodes. Average travel time
falls from 51.1\,s to 46.1\,s (9.8\% reduction). This confirms a clear learning
signal: the MA2C agent improves its policy with experience, answering RQ2
affirmatively.

\begin{table}[h]
\caption{MA2C Training Metrics per Episode}
\label{tab:training}
\centering
\renewcommand{\arraystretch}{1.2}
\begin{tabular}{ccc}
\toprule
\textbf{Episode} & \textbf{Avg Wait (s)} & \textbf{Avg Travel (s)} \\
\midrule
1  & 9.87 & 51.06 \\
2  & 9.51 & 50.77 \\
4  & 8.27 & 48.90 \\
6  & 6.82 & 47.05 \\
8  & 7.07 & 47.36 \\
10 & 7.89 & 48.46 \\
12 & 7.94 & 48.60 \\
14 & \textbf{6.19} & \textbf{46.13} \\
\bottomrule
\end{tabular}
\end{table}

\subsection{RQ1: Comparison with Greedy Baseline}

Table~\ref{tab:comparison} compares the deterministic MA2C evaluation against
the greedy baseline. The greedy controller achieves lower absolute waiting and
travel times than the partially-trained MA2C agent.

\begin{table}[h]
\caption{MA2C vs.\ Greedy Baseline}
\label{tab:comparison}
\centering
\renewcommand{\arraystretch}{1.2}
\begin{tabular}{lcc}
\toprule
\textbf{Controller} & \textbf{Avg Wait (s)} & \textbf{Avg Travel (s)} \\
\midrule
Greedy (2 seeds)     & 0.71  & 38.98 \\
MA2C eval (10 seeds) & 6.06  & 49.03 \\
\bottomrule
\end{tabular}
\end{table}

These results reflect the limited training budget of 14 episodes, substantially
below the 694 episodes used by Chu~\textit{et~al.}~\cite{chu2020} on the Sioux
Falls benchmark. The computational cost of simulating a full-scale real-world
road network (442 detectors, 190 agents, 3600-second episodes) on consumer
hardware constrains the number of training episodes achievable within the
prototype scope. Importantly, the training trend in Table~\ref{tab:training} is
still descending at episode~14, indicating the agent has not yet converged and
would continue to improve with further training.

It is also worth noting that the two metrics are not measured identically: the
greedy waiting time is computed as the ratio of step-level halting vehicle counts
to total departures, while the MA2C evaluation figure comes from SUMO's
per-vehicle \texttt{tripinfo} tracker. Both metrics favour the greedy controller,
so the direction of comparison is sound, but exact magnitudes should be
interpreted accordingly.

\subsection{Comparison with State of the Art}

Du~\textit{et~al.}~\cite{du2023} report 14--54\% reductions in congestion
severity for GAQ-EBkSP over rule-based baselines after full convergence. The
37.3\% reduction in waiting time observed here over just 14 episodes is
consistent in direction with those findings, though direct numerical comparison
is not possible given the different networks, demand patterns, and training
scales.

%% -----------------------------------------------------------
%%  V. DISCUSSION
%% -----------------------------------------------------------
\section{Discussion}

The results position TORRS as a proof-of-concept demonstration that the MA2C
framework is applicable to a real-world island road network. The 37.3\%
reduction in waiting time across 14 episodes establishes a clear learning
signal, and the still-descending training curve at episode~14 indicates that
substantially greater improvements would be achievable given adequate
computational resources.

The greedy baseline outperforms the 14-episode MA2C prototype in absolute terms,
which is expected at this stage of training. This is consistent with the broader
RL literature: partially trained RL agents rarely surpass well-designed
heuristics, but converged RL agents typically outperform them significantly. The
Chu~\textit{et~al.} reference implementation required approximately 694 episodes
to converge on a simpler synthetic network, compared with the 14 episodes
completed here on Malta's full-scale real-world topology.

The compliance sensitivity analysis (varying the proportion of vehicles following
routing recommendations) and seasonal demand variation remain important
directions for future work. If TORRS can demonstrate significant network
improvement at low compliance levels (e.g., 20--40\%), it would provide Transport
Malta with a strong evidence base for piloting connected-vehicle routing guidance
without requiring island-wide adoption. The explicit secondary-bottleneck design
of the MA2C reward function, motivated by Gregucci~\textit{et~al.}~\cite{gregucci2023},
is expected to become increasingly valuable as training progresses and the agent
moves from random exploration toward deliberate policy exploitation.

Unlike wide-grid urban networks with many substitute corridors, Malta's bounded
topology amplifies the risk that improperly designed rerouting displaces a
bottleneck rather than dissolving it. This makes network-wide cooperative
learning, as provided by MA2C's fingerprinting and shared reward structure,
particularly important in the Malta context.

%% -----------------------------------------------------------
%%  VI. CONCLUSION
%% -----------------------------------------------------------
\section{Conclusion}

This paper has presented the design, implementation, and preliminary evaluation
of TORRS, a Multi-Agent Advantage Actor-Critic system for joint traffic signal
control and vehicle routing on Malta's real road network. The key contributions
are: a working SUMO microsimulation of Malta imported from OpenStreetMap; a
successful adaptation of the MA2C algorithm to 190 cooperative agents on a
full-scale island road network; and an empirical demonstration of learning, with
a 37.3\% reduction in average vehicle waiting time over 14 training episodes.

The primary limitation is computational: simulating a full-scale island road
network at 5-second resolution requires approximately 30--40 minutes of
wall-clock time per episode on consumer CPU hardware, restricting the prototype
to 14 training episodes. Full convergence, as achieved in the original MA2C
paper~\cite{chu2020}, requires on the order of 500--700 episodes, representing
several weeks of continuous computation under the current setup.

Future work should address three directions. First, training should be extended
using cloud-based GPU infrastructure to achieve policy convergence and establish
a definitive comparison with the greedy baseline. Second, seasonal demand
variation should be modelled using separate base-season and peak-tourist-season
OD matrices to test the robustness of the learned policy to demand surges.
Third, a compliance sensitivity analysis should be completed by varying the
proportion of vehicles following routing recommendations from 10\% to 100\%,
to provide actionable guidance on the minimum adoption threshold required for
measurable network-wide benefit.

%% -----------------------------------------------------------
%%  REFERENCES
%% -----------------------------------------------------------
\begin{thebibliography}{99}

\bibitem{tomtom2026}
TomTom,
``Traffic Index: Malta Country Report 2026,''
[Online]. Available: \url{https://www.tomtom.com/traffic-index/country/malta}.
[Accessed: 11-May-2026].

\bibitem{nso2024}
National Statistics Office Malta,
``Vehicle registrations and traffic statistics,''
NSO Malta, Valletta, Malta, 2024.

\bibitem{zulkarnain2021}
M. Zulkarnain, M. A. Mohamed, and S. A. Sulaiman,
``Intelligent Transportation Systems (ITS): A systematic review using a
large-scale NLP framework,''
\textit{PLOS ONE}, vol.~16, no.~12, p.~e0261053, Dec. 2021.
[Online]. Available: \url{https://pmc.ncbi.nlm.nih.gov/articles/PMC8695271}

\bibitem{du2023}
B. Du, Y. Liu, Y. Xu, X. Luo, and M. Kowalczyk,
``Dynamic urban traffic rerouting with fog-cloud reinforcement learning,''
\textit{Computer-Aided Civil and Infrastructure Engineering},
vol.~38, no.~12, pp.~1607--1626, 2023.
DOI: \url{https://doi.org/10.1111/mice.13115}

\bibitem{gregucci2023}
D. Gregucci, L. Bedogni, L. Bononi, and M. Di~Felice,
``Dynamic adaptive vehicle re-routing strategy for traffic congestion
mitigation,''
\textit{Transportation Research Interdisciplinary Perspectives},
vol.~18, p.~100786, 2023.
DOI: \url{https://doi.org/10.1016/j.trip.2023.100786}

\bibitem{chu2020}
T. Chu, J. Wang, L. Codec\`{a}, and Z. Li,
``Multi-agent deep reinforcement learning for large-scale traffic signal
control,''
\textit{IEEE Transactions on Intelligent Transportation Systems},
vol.~21, no.~3, pp.~1086--1095, Mar. 2020.
DOI: \url{https://doi.org/10.1109/TITS.2019.2901791}

\bibitem{lopez2018}
P.~A. Lopez \textit{et al.},
``Microscopic traffic simulation using SUMO,''
in \textit{Proc. IEEE Int. Conf. Intelligent Transportation Systems (ITSC)},
Maui, HI, USA, 2018, pp.~2575--2582.

\end{thebibliography}

\end{document}
